% Generated draft: requires mathematical and source review.
\documentclass[11pt]{article}
\usepackage[margin=25mm]{geometry}
\usepackage{fontspec}
\setmainfont{Latin Modern Roman}
\usepackage{amsmath,amssymb,mathtools,unicode-math}
\setmathfont{Latin Modern Math}
\usepackage{enumitem,needspace,xurl,hyperref}
\setlength{\parindent}{0pt}
\setlength{\parskip}{5pt}
\setlength{\emergencystretch}{3em}
\newcommand{\N}{\mathbb{N}}
\newcommand{\Z}{\mathbb{Z}}
\newcommand{\Q}{\mathbb{Q}}
\newcommand{\R}{\mathbb{R}}
\newcommand{\C}{\mathbb{C}}
\newcommand{\F}{\mathbb{F}}
\newcommand{\A}{\mathbb{A}}
\newcommand{\PP}{\mathbb{P}}
\newcommand{\E}{\mathbb{E}}
\newcommand{\eps}{\varepsilon}
\newcommand{\dd}{\,\mathrm{d}}
\begin{document}
\section*{UnsolvedMath ID: 701547}
\textbf{Source problem code:} catalogue entry
\subsection{Definitions and mathematical statement}
For elements of a group, use the convention \([x,y]=x^{-1}y^{-1}xy\). A multilinear commutator word of weight \(n\) is a fully bracketed commutator expression in \(n\) distinct variables, each occurring exactly once. Equivalently, such expressions are built recursively by taking commutators \([u,v]\) of expressions involving disjoint sets of variables, starting with individual variables. Allowing an individual variable as a word of weight one gives a trivial instance of the question.

For a group \(G\) and such a word \(w=w(x_1,\ldots,x_n)\), define the set of \(w\)-values and its verbal subgroup by
\[
G_w=\{w(g_1,\ldots,g_n):g_1,\ldots,g_n\in G\},
\qquad
w(G)=\langle G_w\rangle.
\]
For a finite group \(H\) and a prime \(p\), a Sylow \(p\)-subgroup is a subgroup of order equal to the largest power of \(p\) dividing \(|H|\).

\textbf{Question (P. Shumyatsky, Kourovka Notebook Problem 18.119).} For every multilinear commutator word \(w\), every finite group \(G\), every prime \(p\), and every Sylow \(p\)-subgroup \(P\) of \(w(G)\), must
\[
P=\bigl\langle P\cap G_w\bigr\rangle
\]
hold? Here the generating values are evaluated on elements of the ambient group \(G\) and must themselves lie in \(P\); the arguments of these evaluations are not required to lie in \(P\). If \(p\nmid |w(G)|\), then \(P=1\) and the assertion is automatic.
\subsection{Short English statement}
Given a finite group and a multilinear commutator word, must every Sylow subgroup of the associated verbal subgroup be generated by word values that belong to that Sylow subgroup? The evaluations may use any elements of the ambient group. The supplied registry labels the problem open; this draft does not independently verify that status.
\subsection{Sources}
\begin{itemize}
\item Supplied local UnsolvedMath registry record, permanent ID \(701547\), exact source code \texttt{catalogue entry}, title \emph{Kourovka Notebook Problem 18.119}. The record attributes the problem to P. Shumyatsky, gives the proposed year as 2014, and labels its status \texttt{open}.
\item The supplied background cites \emph{The Kourovka Notebook (Unsolved Problems in Group Theory)}, 18th issue (2014), problem 18.119, and the 21st issue edition (2026), identified as \texttt{arXiv:1401.0300v48}. The supplied collection description names E. I. Khukhro and V. D. Mazurov as editors. These bibliographic details are supplied attribution and have not been independently verified.
\end{itemize}
\end{document}
